\chapter{Marco Teórico}

\section{Internet de las Cosas (IoT)}
El Internet de las Cosas se refiere a la interconexión digital de objetos cotidianos con internet. En este proyecto, la estación meteorológica actúa como un dispositivo "edge" que procesa y transmite datos hacia una arquitectura centralizada en la nube.

\section{Hardware y Sensores}

\subsection{ESP32}
El ESP32 es un microcontrolador de bajo costo y bajo consumo de energía con Wi-Fi y Bluetooth de modo dual integrado. Es el corazón del proyecto debido a su potencia de procesamiento y soporte nativo para protocolos de red.

\subsection{BME280: Sensor Ambiental}
El sensor BME280 de Bosch es un sensor digital integrado que mide presión, humedad y temperatura. Utiliza el protocolo I2C para la comunicación con el microcontrolador.

\subsection{MPU6500: Acelerómetro y Giroscopio}
Este sensor permite detectar la orientación y vibraciones de la estación, lo cual es útil para verificar la estabilidad física de la estructura ante ráfagas de viento.

\subsection{INMP441: Micrófono Digital I2S}
A diferencia de los micrófonos analógicos, el INMP441 utiliza el protocolo I2S (Inter-IC Sound), lo que permite una captura de audio digital de alta fidelidad directamente en el microcontrolador sin necesidad de conversores ADC externos.

\section{Infraestructura de Datos}

\subsection{Supabase y PostgreSQL}
Supabase es una alternativa de código abierto a Firebase que utiliza PostgreSQL como base de datos central. Permite capacidades de "Realtime" que son cruciales para actualizar el dashboard en el instante en que el sensor realiza una lectura.

\subsection{WebSockets}
Los WebSockets permiten una comunicación bidireccional permanente entre el cliente (Dashboard) y el servidor, eliminando la necesidad de realizar "polling" constante y reduciendo la latencia de la aplicación.
